% Document A (readable edition). Lean counterpart: lean/GL2TF/VolI/P0/Conventions.lean
% Chapter files of A and B have identical names; owner per sprint: see the sprint note.
% § "Local maximal compact subgroups" written (M1, unit U1); M2 (unit U1): B, T, N, w, K_inf, k_theta; M3 (unit U1): the rest.

\chapter{Conventions and the algebraic groups}\label{ch:ch-0-01-conventions}

\noindent\emph{Sources:} \cite[§§2, 5]{AIntro}; \cite[Lecture IV §1]{Gel}.

\section{\texorpdfstring{The group $\GL_2$ and its standard subgroups}{The group GL₂ and its standard subgroups}}\label{sec:conv-subgroups}

\noindent\emph{Sources for this section:} \cite[§2]{AIntro}; \cite[Lecture IV §1]{Gel}; \cite[§5]{JL} (the
archimedean notation $n(x)$, $\diag(a_1,a_2)$, $w$).

Throughout, $R$ is a commutative ring and $\GL_2(R)=\Mat_2(R)^\times$; a matrix $g\in\Mat_2(R)$ lies in $\GL_2(R)$ if
and only if $\det g\in R^\times$ (the inverse is $(\det g)^{-1}\operatorname{adj}(g)$). The definitions below are
used for $R=\R$ in milestone M2 and for local and adelic rings later.

\begin{definition}[the group $\GL_2$]
\tfstmt{I.conv.def:gl2}
\end{definition}
\noindent\emph{Homomorphism.} Applying $f$ entrywise to a matrix commutes with sums and products of matrices and
with determinants, because $f$ is a ring homomorphism; it maps $1$ to $1$. So $f$ maps an invertible matrix $g$ with inverse
$h$ to a matrix $f_*g$ with $f_*g\cdot f_*h=f_*(gh)=1=f_*h\cdot f_*g$, and $f_*(gg')=f_*g\,f_*g'$. Moreover
$(\mathrm{id}_R)_*=\mathrm{id}$ and $(f'\circ f)_*=f'_*\circ f_*$ (entrywise), so $R\mapsto\GL_2(R)$ is a functor. In
the library this is \texttt{Matrix.GeneralLinearGroup.map}; the group $\GL_2(R)$ itself is \texttt{GL (Fin 2) R}, with no
project alias.

\begin{definition}[the Borel subgroup]
\tfstmt{I.conv.def:borel}
\end{definition}
\noindent\emph{Subgroup.} $1\in B(R)$. For $g,h$ with $g_{21}=h_{21}=0$, $(gh)_{21}=g_{21}h_{11}+g_{22}h_{21}=0$, and
$(g^{-1})_{21}=-(\det g)^{-1}g_{21}=0$.

\begin{definition}[the diagonal torus]
\tfstmt{I.conv.def:torus}
\end{definition}
\noindent\emph{Subgroup.} Products and inverses of invertible diagonal matrices are diagonal; a diagonal matrix
$\diag(t_1,t_2)$ is invertible if and only if $t_1t_2\in R^\times$, i.e.\ $t_1,t_2\in R^\times$.

\begin{definition}[the unipotent subgroup]
\tfstmt{I.conv.def:unipotent}
\end{definition}
\noindent\emph{Proof of the claims.} $n(x)n(y)=n(x+y)$ by multiplying out, and $n(0)=1$; hence $n(x)$ is invertible
with inverse $n(-x)$, $n$ is a homomorphism, and its image $N(R)$ is a subgroup. $n$ is injective because $x$ is the
$(1,2)$-entry of $n(x)$.

\begin{definition}[the Weyl element]
\tfstmt{I.conv.def:weyl-element}
\end{definition}
\noindent $\det w=0\cdot0-1\cdot(-1)=1$, so $w\in\GL_2(R)$; also $w^2=-1$ and $w^{-1}=-w$.

\begin{definition}[the centre]
\tfstmt{I.conv.def:center}
\uses{I.conv.def:gl2}
\end{definition}
\noindent\emph{Subgroup.} $(z\cdot1)(z'\cdot1)=zz'\cdot1$ and $(z\cdot1)^{-1}=z^{-1}\cdot1$. \emph{It is the centre:}
scalar matrices commute with every matrix. Conversely, if $g=\begin{psmallmatrix}a&b\\c&d\end{psmallmatrix}$ commutes
with $n(1)$, comparing $gn(1)=\begin{psmallmatrix}a&a+b\\c&c+d\end{psmallmatrix}$ with
$n(1)g=\begin{psmallmatrix}a+c&b+d\\c&d\end{psmallmatrix}$ gives $c=0$ and $a=d$; commuting with $\begin{psmallmatrix}1&0\\1&1\end{psmallmatrix}$ gives
$b=0$ in the same way. So $g=a\cdot1$ with $a^2=\det g\in R^\times$, hence $a\in R^\times$. (This remark is not used
formally.)

\begin{definition}[the opposite Borel subgroup]
\tfstmt{I.conv.def:opposite-borel}
\uses{I.conv.def:gl2}
\end{definition}
\noindent\emph{Proof of the claims.} $\Bbar(R)$ is the image of $B(R)$ under the transpose anti-automorphism $g\mapsto g^{\mathsf T}$ of
$\GL_2(R)$ (it reverses products and preserves inverses), so it is a subgroup by \cref{I.conv.def:borel}. As for
$n$, $\bar n(x)\bar n(y)=\bar n(x+y)$ and $\bar n(0)=1$, and $\bar n$ is injective since $x$ is the $(2,1)$-entry
of $\bar n(x)$.

\begin{lemma}[functoriality of the subgroups]
\tfstmt{I.conv.lem:subgroups-functorial}
\uses{I.conv.def:gl2, I.conv.def:borel, I.conv.def:torus, I.conv.def:unipotent, I.conv.def:center, I.conv.def:opposite-borel}
\end{lemma}
\begin{proof}
\uses{I.conv.def:gl2, I.conv.def:borel, I.conv.def:torus, I.conv.def:unipotent, I.conv.def:center, I.conv.def:opposite-borel}
Since $f(0)=0$ and $f_*$ acts entrywise, a vanishing entry of $g$ gives a vanishing entry of $f_*g$: the conditions
$g_{21}=0$ (for $B$), $g_{12}=g_{21}=0$ (for $T$) and $g_{12}=0$ (for $\Bbar$) are preserved. The formulas follow
from $f(0)=0$, $f(1)=1$ by computing entrywise; they show that $N$, $\Nbar$ and $T$ are preserved
($T(R)$ consists of the $\diag(t_1,t_2)$ by \cref{I.conv.def:torus}). Finally
$f_*(z\cdot1)=f(z)\cdot1$, and $f(z)$ is a unit with inverse $f(z^{-1})$.
\end{proof}

\section{\texorpdfstring{Structure of $B$ and of $N_G(T)$}{Structure of B and of NG(T)}}\label{sec:conv-structure}

\noindent\emph{Sources for this section:} \cite[§2]{AIntro}; \cite[Lecture IV §1]{Gel}. The arguments are explicit
$2\times2$ computations, valid over every commutative ring unless a field is required.

\begin{lemma}[$T$ normalizes $N$]
\tfstmt{I.conv.lem:torus-normalizes-unipotent}
\uses{I.conv.def:torus, I.conv.def:unipotent}
\end{lemma}
\begin{proof}
\uses{I.conv.def:torus, I.conv.def:unipotent}
Write $t=\diag(t_1,t_2)$, so $t^{-1}=\diag(t_1^{-1},t_2^{-1})$. Multiplying out,
$t\,n(x)=\begin{psmallmatrix}t_1&t_1x\\0&t_2\end{psmallmatrix}$ and
$t\,n(x)\,t^{-1}=\begin{psmallmatrix}1&t_1xt_2^{-1}\\0&1\end{psmallmatrix}=n(t_1t_2^{-1}x)$. Hence
$tN(R)t^{-1}\subseteq N(R)$ for every $t\in T(R)$; applying this to $t^{-1}$ gives $t^{-1}N(R)t\subseteq N(R)$, i.e.
$N(R)\subseteq tN(R)t^{-1}$.
\end{proof}

\begin{theorem}[$B=T\ltimes N$]
\tfstmt{I.conv.thm:borel-semidirect}
\uses{I.conv.def:borel, I.conv.def:torus, I.conv.def:unipotent}
\end{theorem}
\begin{proof}
\uses{I.conv.lem:torus-normalizes-unipotent, I.conv.def:borel, I.conv.def:torus, I.conv.def:unipotent}
(ii) An element of $T(R)\cap N(R)$ is diagonal and of the form $n(x)$; its $(1,2)$-entry is $x=0$, so it is $n(0)=1$.

(i) Let $b=\begin{psmallmatrix}a&y\\0&d\end{psmallmatrix}\in B(R)$. Then $\det b=ad\in R^\times$, so $a$ and $d$ are
units (with inverses $d(ad)^{-1}$ and $a(ad)^{-1}$). Multiplying out,
$\diag(a,d)\,n(a^{-1}y)=\begin{psmallmatrix}a&y\\0&d\end{psmallmatrix}=b$. Uniqueness: if $tn=t'n'$ then
$t'^{-1}t=n'n^{-1}\in T(R)\cap N(R)=\{1\}$ by (ii), so $t=t'$ and $n=n'$.

(iii) $N(R)\subseteq B(R)$ since $n(x)_{21}=0$. For $b=tn_0\in B(R)$ as in (i) and $n\in N(R)$,
$bnb^{-1}=t(n_0nn_0^{-1})t^{-1}\in tN(R)t^{-1}=N(R)$ by \cref{I.conv.lem:torus-normalizes-unipotent}, since
$N(R)$ is a subgroup.
\end{proof}
\noindent By (i)--(iii), $B(R)=T(R)\ltimes N(R)$; also $b=n(yd^{-1})\diag(a,d)$, so every $b$ is uniquely $nt$ as
well (apply (i) to $b^{-1}$ and invert). The identification with the semidirect product $N\rtimes T$ as a group isomorphism
is plumbing (B only).

\begin{lemma}[$B\cap\Nbar=1$]
\tfstmt{I.conv.lem:borel-inter-opposite-unipotent}
\uses{I.conv.def:borel, I.conv.def:opposite-borel}
\end{lemma}
\begin{proof}
\uses{I.conv.def:borel, I.conv.def:opposite-borel}
If $\bar n(x)\in B(R)$, its $(2,1)$-entry $x$ vanishes, so $\bar n(x)=\bar n(0)=1$.
\end{proof}

\begin{theorem}[the normalizer of the torus]
\tfstmt{I.conv.thm:normalizer-torus}
\uses{I.conv.def:torus, I.conv.def:weyl-element}
\end{theorem}
\begin{proof}
\uses{I.conv.def:torus, I.conv.def:weyl-element}
\emph{$T$ and $w$ normalize $T$.} $T(k)$ is a subgroup. For $t=\diag(t_1,t_2)$, with
$w^{-1}=\begin{psmallmatrix}0&-1\\1&0\end{psmallmatrix}$ one computes $wtw^{-1}=\diag(t_2,t_1)\in T(k)$, so
$wT(k)w^{-1}\subseteq T(k)$, and equality holds because $w^{-1}tw=\diag(t_2,t_1)$ as well. Hence every element of
$T(k)\cup wT(k)$ normalizes $T(k)$.

\emph{Conversely,} let $g$ normalize $T(k)$. Since $k$ has at least three elements there is $a\in k$ with
$a\neq0,1$. Put $t_0=\diag(a,1)$ and $s=gt_0g^{-1}\in T(k)$, say $s=\diag(b_1,b_2)$. Conjugate matrices have the same
trace and determinant: $b_1+b_2=a+1$ and $b_1b_2=a$. Hence $(X-b_1)(X-b_2)=X^2-(a+1)X+a=(X-a)(X-1)$, and since $k$ is
a field, $(b_1,b_2)=(a,1)$ or $(b_1,b_2)=(1,a)$.
\begin{itemize}
\item If $s=t_0$, then $g$ commutes with $t_0$. For $g=\begin{psmallmatrix}\alpha&\beta\\\gamma&\delta\end{psmallmatrix}$,
  $gt_0=\begin{psmallmatrix}a\alpha&\beta\\a\gamma&\delta\end{psmallmatrix}$ and
  $t_0g=\begin{psmallmatrix}a\alpha&a\beta\\\gamma&\delta\end{psmallmatrix}$, so $(a-1)\beta=(a-1)\gamma=0$ and
  $\beta=\gamma=0$ as $a\neq1$: $g\in T(k)$.
\item If $s=\diag(1,a)=wt_0w^{-1}$, then $h=w^{-1}g$ satisfies $ht_0h^{-1}=w^{-1}sw=t_0$, so $h\in T(k)$ by the first
  case, and $g=wh\in wT(k)$.
\end{itemize}
\emph{Disjointness.} Elements of $wT(k)$ are $w\diag(t_1,t_2)=\begin{psmallmatrix}0&t_2\\-t_1&0\end{psmallmatrix}$,
whose $(1,2)$-entry $t_2$ is non-zero, while elements of $T(k)$ have $(1,2)$-entry $0$.
\end{proof}
\noindent Over $\mathbb F_2$ the torus is trivial and its normalizer is all of $\GL_2(\mathbb F_2)$, so the hypothesis
on $k$ cannot be dropped.

\section{Local maximal compact subgroups}\label{sec:conv-max-compact}

\noindent\emph{Sources for this section:} \cite[§9.1]{Cas-p}; \cite[§9]{BumpH}.

The non-archimedean maximal compact subgroup is defined for an arbitrary ring homomorphism $R\to F$, so that this
chapter does not depend on the theory of local fields (\cref{ch:ch-0-02-local-field}); the case of interest is
$\Ocal_F\subseteq F$ for a non-archimedean local field $F$.

\begin{definition}[the subgroup $K(R)$]
\tfstmt{I.conv.def:max-compact-local}
\end{definition}
\noindent\emph{Notation.} For a non-archimedean local field $F$ (\cref{I.lf.def:nalf}) we write
$K_F=K(\Ocal_F)$, for the inclusion $\Ocal_F\subseteq F$; this is a notation for an instance of the definition, not
a separate definition. For $F=\Qp$ the ring of integers is $\Zp$, and $K_p=K_{\Qp}=\GL_2(\Zp)$ is the group of the notation table.

\begin{lemma}[membership in $K(R)$]
\tfstmt{I.conv.lem:max-compact-local-mem}
\uses{I.conv.def:max-compact-local}
\end{lemma}
\begin{proof}
\uses{I.conv.def:max-compact-local}
For a matrix $A$ over $R$ write $\iota(A)$ for the matrix over $F$ obtained by applying $\iota$ entrywise; then
$\iota(AB)=\iota(A)\iota(B)$, $\iota(1)=1$ and $\det\iota(A)=\iota(\det A)$, and $A\mapsto\iota(A)$ is injective.

(i)$\Rightarrow$(iii): if $g=\iota(A)$ with $A\in\GL_2(R)$, then $g^{-1}=\iota(A^{-1})$, and both have entries
in $\iota(R)$.

(iii)$\Rightarrow$(ii): write $g=\iota(A)$ and $g^{-1}=\iota(B)$ with $A,B\in\Mat_2(R)$. Then
$\iota(AB)=gg^{-1}=1=\iota(1)$, hence $AB=1$ by injectivity, so $\det A\cdot\det B=1$ and $\det A\in R^\times$. Thus
$\det g=\iota(\det A)\in\iota(R^\times)$.

(ii)$\Rightarrow$(i): write $g=\iota(A)$ with $A\in\Mat_2(R)$ and $\det g=\iota(r)$ with $r\in R^\times$. Then
$\iota(\det A)=\det g=\iota(r)$, so $\det A=r$ is a unit, and $A$ is invertible in $\Mat_2(R)$ with inverse
$r^{-1}\operatorname{adj}(A)$. Hence $A\in\GL_2(R)$ and $g=\iota(A)\in K(R)$.
\end{proof}
\noindent In particular, for a non-archimedean local field $F$: $g\in K_F$ if and only if $g\in\Mat_2(\Ocal_F)$ and
$\det g\in\Ocal_F^\times$. Consequently $K_F$ contains the permutation matrix
$s=\begin{psmallmatrix}0&1\\1&0\end{psmallmatrix}$, the elementary matrices
$\begin{psmallmatrix}1&x\\0&1\end{psmallmatrix}$, $\begin{psmallmatrix}1&0\\x&1\end{psmallmatrix}$ ($x\in\Ocal_F$),
and the diagonal matrices $\diag(u_1,u_2)$ ($u_i\in\Ocal_F^\times$). It also contains the scalar matrices $u\cdot1$
with $u\in\Ocal_F^\times$.

\subsection*{The archimedean maximal compact subgroup}

\noindent\emph{Sources:} \cite[§5]{JL} (the orthogonal group $\Ort(2,\R)$ and the characters
$\begin{psmallmatrix}\cos\theta&\sin\theta\\-\sin\theta&\cos\theta\end{psmallmatrix}\mapsto e^{in\theta}$); \cite[§5.2]{GH}
(the same rotation $\kappa_\theta$). The sign convention of $k_\theta$ is that of both sources; with it, $k_\theta=\exp(\theta U)$
for $U=\begin{psmallmatrix}0&1\\-1&0\end{psmallmatrix}$, and the weight characters (of the parent's Vol.~I, Ch.~1.4; not used here) are
$k_\theta\mapsto e^{in\theta}$.

\begin{definition}[$\Kinf=\Ort(2)$]
\tfstmt{I.conv.def:max-compact-real}
\end{definition}
\noindent\emph{Subgroups.} For square matrices $k^{\mathsf T}k=1$ implies $kk^{\mathsf T}=1$, so $k^{-1}=k^{\mathsf T}$ and
$(k^{-1})^{\mathsf T}k^{-1}=kk^{\mathsf T}=1$; and $(k_1k_2)^{\mathsf T}k_1k_2=k_2^{\mathsf T}(k_1^{\mathsf T}k_1)k_2=1$. $\mathrm{SO}(2)$ is the
kernel of the homomorphism $\det|_{\Kinf}$.

\begin{definition}[rotations]
\tfstmt{I.conv.def:rotation}
\uses{I.conv.def:max-compact-real}
\end{definition}

\begin{remark}[structure of $\Ort(2)$]\label{rem:rotation}
Put $\varepsilon=\diag(1,-1)$ and $U=\begin{psmallmatrix}0&1\\-1&0\end{psmallmatrix}$.
\begin{enumerate}[label=(\alph*)]
\item\label{rot-hom} $k_\theta\in\mathrm{SO}(2)$, $k_0=1$ and $k_\theta k_\varphi=k_{\theta+\varphi}$; so $\theta\mapsto k_\theta$ is a
  homomorphism $\R\to\mathrm{SO}(2)$, continuous (its entries are).
\item\label{rot-shape} $\mathrm{SO}(2)=\bigl\{\begin{psmallmatrix}a&b\\-b&a\end{psmallmatrix}:a^2+b^2=1\bigr\}$; in particular $k\in\mathrm{SO}(2)$ is
  determined by its bottom row.
\item\label{rot-surj} Every $k\in\mathrm{SO}(2)$ is $k_\theta$ for some $\theta\in\R$, and $k_\theta=1$ if and only if $\theta\in2\pi\Z$.
\item\label{rot-o2} $\det k=\pm1$ for $k\in\Kinf$; every $k\in\Kinf$ is $k_\theta$ or $k_\theta\varepsilon$ for some $\theta$;
  and $\varepsilon k_\theta\varepsilon=k_{-\theta}$.
\item\label{rot-exp} $k_\theta=\exp(\theta U)$.
\end{enumerate}
\emph{Proof.} (a) $\det k_\theta=\cos^2\theta+\sin^2\theta=1$ and $k_\theta^{\mathsf T}k_\theta=1$ by the same identity;
the product formula is the addition theorem for $\cos$ and $\sin$. (b) For $k=\begin{psmallmatrix}a&b\\c&d\end{psmallmatrix}\in\mathrm{SO}(2)$,
$k^{-1}=\begin{psmallmatrix}d&-b\\-c&a\end{psmallmatrix}$ (as $\det k=1$) equals $k^{\mathsf T}=\begin{psmallmatrix}a&c\\b&d\end{psmallmatrix}$, so $d=a$,
$c=-b$, and $\det k=a^2+b^2=1$; conversely such matrices lie in $\mathrm{SO}(2)$. (c) Given $a^2+b^2=1$ choose $\theta$ with
$\cos\theta=a$, $\sin\theta=b$ (polar coordinates of the point $(a,b)$ of the unit circle); then $k=k_\theta$ by (b).
$k_\theta=1$ means $\cos\theta=1$, $\sin\theta=0$, i.e.\ $\theta\in2\pi\Z$. (d) $\det(k)^2=\det(k^{\mathsf T}k)=1$. If
$\det k=-1$ then $k\varepsilon\in\mathrm{SO}(2)$, so $k\varepsilon=k_\theta$ by (c) and $k=k_\theta\varepsilon$ ($\varepsilon^2=1$). The last
identity is a direct multiplication. (e) $U^2=-1$, so the exponential series gives
$\exp(\theta U)=\sum_m\frac{(-1)^m\theta^{2m}}{(2m)!}1+\sum_m\frac{(-1)^m\theta^{2m+1}}{(2m+1)!}U=\cos\theta\cdot1+\sin\theta\cdot U=k_\theta$.
\end{remark}

\section{\texorpdfstring{The space $\aB$, roots and chambers}{The space 𝔞B, roots and chambers}}\label{sec:conv-aB}

\noindent\emph{Sources for this section:} \cite[§§2, 5]{AIntro}; \cite[§1]{A-Tr}. For $\GL_2$ we write everything in
coordinates; the coordinate $h=(h_1,h_2)$ corresponds to the homomorphism $X^*(T)\to\R$,
$(t\mapsto t_1^{m_1}t_2^{m_2})\mapsto m_1h_1+m_2h_2$, which identifies $\Hom(X^*(T)_\Q,\R)$ with $\R^2$.

\begin{definition}[the space $\aB$]
\tfstmt{I.conv.def:space-a-b}
\end{definition}

\begin{definition}[the space $\aBG$ and the projection]
\tfstmt{I.conv.def:space-a-b-g}
\uses{I.conv.def:space-a-b}
\end{definition}
\noindent\emph{Properties.} $\aBG$ is the kernel of the linear form $h\mapsto h_1+h_2$, hence a subspace. The value
$h_B^G$ lies in $\aBG$, it equals $h$ for $h\in\aBG$ (then $h_2=-h_1$ and $\frac{h_1-h_2}2=h_1$), and
$h-h_B^G=\frac{h_1+h_2}2(1,1)$ lies in $\mathfrak a_G=\{(x,x)\}$. So $\aB=\mathfrak a_G\oplus\aBG$ and $h\mapsto h_B^G$
is the projection along $\mathfrak a_G$.

\begin{definition}[the simple root and $\rho$]
\tfstmt{I.conv.def:simple-root}
\uses{I.conv.def:space-a-b}
\end{definition}
\noindent\emph{Consistency.} $\alpha$ is the differential of the root $t\mapsto t_1/t_2$ of $T$ in $\mathrm{Lie}(N)$ (\cref{I.conv.lem:torus-normalizes-unipotent}),
and $\alpha$ vanishes on $\mathfrak a_G$, so $\alpha(h_B^G)=\alpha(h)$. This agrees with the coordinate of decision D15 in
the notation table: $\lambda\in\C$ corresponds to $\lambda\alpha$, and $\rho=\frac12\alpha$ corresponds to $\lambda=\frac12$.
With the height map of \cref{ch:ch-0-05-decompositions}, $e^{2\rho(\HB(t))}=\abs{t_1/t_2}_\A=\deltaB(t)$.

\begin{definition}[the chamber functions]
\tfstmt{I.conv.def:chamber-functions}
\uses{I.conv.def:simple-root, I.conv.def:space-a-b}
\end{definition}
\noindent\emph{Remarks.} Since $\varpi=\frac12\alpha$ is a positive multiple of $\alpha$, $\hat\tau_B=\tau_B$ for $\GL_2$; the two
names are kept because Arthur's identities \cite[§§5--8]{AIntro} distinguish them. Both depend only on $h_B^G$. The
coroot is $\alpha^\vee=(1,-1)\in\aBG$, with $\alpha(\alpha^\vee)=2$ and $\varpi(\alpha^\vee)=1$, which characterizes
$\varpi$ as the fundamental weight on $\aBG$.

\section{Notation table}
\input{notation}
